%% ============================================================================
%% Abstract and keywords (independent file, included inside \begin{frontmatter})
%% ============================================================================

\begin{abstract}
Sparse-view cone-beam CT (CBCT) reconstruction aims to recover a tomographic
volume from an extremely limited set of projection views, which is an
ill-posed inverse problem that traditional methods make it difficult to 
achieve effective reconstruction, and is typically solved with deterministic
deep networks. Probabilistic formulations can in principle quantify
reconstruction uncertainty and exploit dense observations during training,
yet a principled and well-diagnosed probabilistic design for
Gaussian-splatting-based CBCT reconstruction remains missing. In this work
we propose a probabilistic extension of the 3D-Gaussian-splatting CBCT
reconstructor built on a neural-process-style formulation with a tractable
global latent variable. Our training objective combines
(i)~a variational posterior that is conditioned on dense views during
training, (ii)~a sparse-view prior that is used exclusively at deployment,
and (iii)~an auxiliary prior-reconstruction term that supervises the
deployment path directly. We further introduce a dense-latent
neural-process (dense-NP) backbone together with a lightweight
GT-conditioned CVAE variant (P5) that encodes the ground-truth volume into
the variational posterior during training only, so that \emph{no}
ground truth is required at inference time. Systematic diagnostic
instrumentation -- covering KL divergence, cross-case latent statistics,
prior/posterior reconstruction gaps, and zero- and shuffle-perturbation
interventions -- reveals a persistent \emph{posterior-collapse} failure
mode shared by all probabilistic variants under the dense-decoder
conditioning scheme, and identifies the deterministic mean-code control as
the strongest reference. On the LUNA16 CBCT benchmark, our best
GT-conditioned CVAE variant reaches $29.24$\,dB PSNR / $0.863$\,SSIM, within
$0.3$\,dB of the deterministic mean-code reference ($29.54$\,dB /
$0.870$\,SSIM), while providing calibrated deployment-side uncertainty
estimates. Our analysis isolates the exact mechanism behind variational
collapse in splatting-based reconstruction and points to concrete remedies
(context-conditioned decoding, free-bits floors, and deployment-guided
supervision) that are actionable for future probabilistic CBCT models.
\end{abstract}

%% Research highlights
\begin{highlights}
\item A probabilistic Gaussian-splatting framework for sparse-view CBCT
reconstruction with a tractable global latent variable.
\item A GT-conditioned CVAE posterior (P5) that requires no ground truth
or dense views at inference time.
\item Systematic latent diagnostics exposing posterior collapse under
dense-decoder conditioning, plus concrete remedies.
\end{highlights}

%% Keywords
\begin{keyword}
Sparse-view CBCT reconstruction \sep 3D Gaussian splatting \sep
variational inference \sep conditional variational autoencoder \sep
latent neural processes \sep uncertainty estimation
\end{keyword}